% ======================================================================
% fig-ch02-layers.tex
% Chapter 2 (The mental model), the book's ONE Ch 2 figure.
% Concept: what happens between a keystroke and a result. You type
% into a terminal; the shell reads the line, expands it, and finds
% the program on PATH; the kernel starts a process (PID, working
% directory, a copy of the environment, three streams); output
% flows back to the terminal and the exit code back to the shell.
%
% Semantic palette (core three + neutral, per the design system):
%   blue  = the shell, the decision-making main path
%   green = the running process, where the work happens
%   amber = the exit-code report back (status/check)
%   slate = the inert layers (you, terminal, kernel)
%
% Build:  scripts/render.sh fig-ch02-layers.tex
% ======================================================================
\documentclass[tikz,border=8pt]{standalone}
% ======================================================================
% preamble.tex -- shared TikZ design system for the textbook figures.
%
% Usage in a standalone figure:
%   \documentclass[tikz,border=8pt]{standalone}
%   % ======================================================================
% preamble.tex -- shared TikZ design system for the textbook figures.
%
% Usage in a standalone figure:
%   \documentclass[tikz,border=8pt]{standalone}
%   % ======================================================================
% preamble.tex -- shared TikZ design system for the textbook figures.
%
% Usage in a standalone figure:
%   \documentclass[tikz,border=8pt]{standalone}
%   \input{preamble.tex}      % path relative to the figure, or via TEXINPUTS
%   \begin{document}\begin{tikzpicture} ... \end{tikzpicture}\end{document}
%
% Design intent (read this before adding one-off styles):
%   * One palette, used SEMANTICALLY. A colour means a role, never decoration.
%   * Rounded "card" containers, thin borders, soft fills, one sans-serif.
%   * Real Lucide icons (built to PDF by icons/build_icons.py), one family.
%   * Manual layout. We choose spacing and alignment; we do not auto-route.
%
% Never hand-pick a colour in a figure. If you need a new role, add it here.
% ======================================================================

\usepackage[T1]{fontenc}
\usepackage[scaled=0.95]{helvet}      % Helvetica-like sans; swap for the
\renewcommand\familydefault{\sfdefault} % book font if the template differs.
\usepackage{graphicx}
\usepackage{amssymb}

% Letterspacing for small-caps section labels. Guarded so compilation
% still works if microtype is unavailable.
\IfFileExists{microtype.sty}{\usepackage[letterspace=140]{microtype}}{}
\providecommand{\textls}[1]{#1}   % fallback if microtype is absent

\usepackage{tikz}
\usetikzlibrary{positioning, arrows.meta, shapes.geometric, calc, fit,
  backgrounds, decorations.pathreplacing}

% ----------------------------------------------------------------------
% 1. SEMANTIC PALETTE  (keep in sync with palette.md and icons INK)
%    Each role has Fill (soft), Line (border/accent), Text (label).
%    Anchored to book-private/docs/visual-style-guide.md (the book exhibit
%    standard): ink #1F2933, muted #64748B, border #CBD5E1, light fill
%    #F8FAFC, process blue #2563EB, evidence green #059669, stop amber
%    #D97706, hard-failure red. The three CORE roles (prim / data / stop)
%    are the guide's blue / green / amber. dsgn, inf, stor are an EXTENDED
%    tier for rich conceptual maps; the guide's default is at most three
%    semantic colours per exhibit, so reach for the core three first.
% ----------------------------------------------------------------------
\definecolor{tbInk}{HTML}{1F2933}      % icons, strokes, default ink
\definecolor{tbMute}{HTML}{64748B}     % captions, section labels
\definecolor{tbLine}{HTML}{94A3B8}     % arrows (secondary stroke)
\definecolor{tbHair}{HTML}{CBD5E1}     % borders, dividers, faint panels
\definecolor{tbPaper}{HTML}{F8FAFC}    % panel / light background

% --- CORE semantic roles (the guide's three colours + neutral) ---
% neutral / slate -- the familiar, the inert, the "given"
\definecolor{tbNeutFill}{HTML}{F1F5F9}
\definecolor{tbNeutLine}{HTML}{CBD5E1}
\definecolor{tbNeutText}{HTML}{334155}
% primary / process blue -- the main path, question, project spine
\definecolor{tbPrimFill}{HTML}{EFF6FF}
\definecolor{tbPrimLine}{HTML}{2563EB}
\definecolor{tbPrimText}{HTML}{1E3A8A}
% data / evidence green -- sources, evidence, validation, checks
\definecolor{tbDataFill}{HTML}{ECFDF5}
\definecolor{tbDataLine}{HTML}{059669}
\definecolor{tbDataText}{HTML}{065F46}
% stop / amber -- stop gates, warnings, revision triggers
\definecolor{tbStopFill}{HTML}{FEF3C7}
\definecolor{tbStopLine}{HTML}{D97706}
\definecolor{tbStopText}{HTML}{78350F}

% --- EXTENDED tier (use sparingly; never break the 3-colour default
%     without a reason, e.g. a wide conceptual map) ---
% design / indigo -- design logic, comparison, identification
\definecolor{tbDsgnFill}{HTML}{EEF2FF}
\definecolor{tbDsgnLine}{HTML}{4F46E5}
\definecolor{tbDsgnText}{HTML}{3730A3}
% inference / teal -- claims, interpretation, warrants
\definecolor{tbInfFill}{HTML}{F0FDFA}
\definecolor{tbInfLine}{HTML}{0D9488}
\definecolor{tbInfText}{HTML}{115E59}
% store / violet -- the durable record, model, source of truth
\definecolor{tbStorFill}{HTML}{F5F3FF}
\definecolor{tbStorLine}{HTML}{7C3AED}
\definecolor{tbStorText}{HTML}{5B21B6}
% fail / red -- genuine hard failures only (the guide reserves red)
\definecolor{tbFailFill}{HTML}{FEF2F2}
\definecolor{tbFailLine}{HTML}{DC2626}
\definecolor{tbFailText}{HTML}{991B1B}
% backward-compatible alias: tbWarn* -> amber stop family
\colorlet{tbWarnFill}{tbStopFill}
\colorlet{tbWarnLine}{tbStopLine}
\colorlet{tbWarnText}{tbStopText}

% ----------------------------------------------------------------------
% 2. ICONS
%    Build first:  cd icons && python3 build_icons.py
%    Then: \tbicon{file-text}  or  \tbicon[5mm]{database}
%    The icon inherits the recoloured ink baked in at build time.
% ----------------------------------------------------------------------
\graphicspath{{icons/cache/}{./icons/cache/}{../icons/cache/}%
  {../../icons/cache/}{./}}
\newcommand{\tbicon}[2][4mm]{%
  \raisebox{-0.12\height}{\includegraphics[height=#1]{#2}}}

% ----------------------------------------------------------------------
% 3. NODE STYLES
% ----------------------------------------------------------------------
\tikzset{
  % spacing defaults; override per figure with node distance=...
  node distance = 9mm and 13mm,
  %
  % -- base card. Use a semantic family with  tbcard=tbData  etc.
  %    #1 family stem: Neut Prim Data Stop (core) | Dsgn Inf Stor Fail (ext).
  %    Border 0.7pt (secondary stroke, per the guide's 0.75-1.0pt band).
  tbcard/.style 2 args = {
    rounded corners = 4pt,
    draw = #1Line, line width = 0.7pt,
    fill = #1Fill, text = #1Text,
    align = #2, inner sep = 6pt,
    minimum height = 11mm, minimum width = 26mm,
    font = \sffamily\small
  },
  tbcard/.default = {tbNeut}{center},
  % convenience: one-argument families, centered text
  neut/.style = {tbcard={tbNeut}{center}},
  prim/.style = {tbcard={tbPrim}{center}},
  data/.style = {tbcard={tbData}{center}},
  stop/.style = {tbcard={tbStop}{center}},
  dsgn/.style = {tbcard={tbDsgn}{center}},
  inf/.style  = {tbcard={tbInf}{center}},
  stor/.style = {tbcard={tbStor}{center}},
  fail/.style = {tbcard={tbFail}{center}},
  warn/.style = {tbcard={tbStop}{center}},   % alias -> stop (amber)
  % left-aligned body variants (for cards with a title + detail lines)
  neutL/.style = {tbcard={tbNeut}{left}},
  primL/.style = {tbcard={tbPrim}{left}},
  dataL/.style = {tbcard={tbData}{left}},
  stopL/.style = {tbcard={tbStop}{left}},
  dsgnL/.style = {tbcard={tbDsgn}{left}},
  infL/.style  = {tbcard={tbInf}{left}},
  storL/.style = {tbcard={tbStor}{left}},
  failL/.style = {tbcard={tbFail}{left}},
  %
  % -- circular variable node for causal DAGs. Single family argument:
  %    tbvar=tbPrim (treatment), tbvar=tbInf (outcome), etc.
  tbvar/.style = {
    circle, draw = #1Line, fill = #1Fill, text = #1Text,
    line width = 0.7pt, align = center, inner sep = 1.5pt,
    minimum size = 15mm, font = \sffamily\small
  },
  tbvar/.default = tbNeut,
  %
  % -- emphasis: a heavier border for the "source of truth" / final node
  tbstrong/.style = {line width = 1.5pt},
  %
  % -- code / monospace inset (e.g. a record or schema snippet)
  tbcode/.style = {
    rounded corners = 3pt, draw = tbHair, line width = 0.5pt,
    fill = white, text = tbNeutText, align = left,
    inner sep = 5pt, font = \ttfamily\scriptsize
  },
  %
  % -- step chip: small numbered circle pinned to a card corner
  tbchip/.style = {
    circle, text = white, font = \sffamily\tiny\bfseries,
    inner sep = 0pt, minimum size = 4.2mm, fill = tbLine
  },
  %
  % -- pill tag (status label like "required" / "gated path")
  tbpill/.style = {
    rounded corners = 5pt, draw = #1Line, line width = 0.5pt,
    fill = #1Fill, text = #1Text, inner xsep = 4pt, inner ysep = 1.5pt,
    font = \sffamily\fontsize{6.5}{7.5}\selectfont
  },
  tbpill/.default = tbNeut,
}

% ----------------------------------------------------------------------
% 4. EDGES
% ----------------------------------------------------------------------
% Line-weight convention follows the guide: main path 1.25-1.5pt, secondary
% 0.75-1.0pt; solid = primary sequence / dependency, dashed = feedback /
% checking / optional. Heavier stroke marks the main path.
\tikzset{
  tbedge/.style = {-{Stealth[length=2.6mm, width=2.1mm]},
    draw = tbLine, line width = 0.9pt},               % secondary stroke
  tbedge thick/.style = {tbedge, line width = 1.4pt}, % main path
  tbdash/.style = {tbedge, dashed},                   % soft / checking link
  tbback/.style = {tbedge, dashed, draw = tbStopLine},% revision / feedback
  tbline/.style = {draw = tbLine, line width = 0.9pt},% undirected tie
}

% ----------------------------------------------------------------------
% 5. SECTIONS, TITLES, FOOTER
% ----------------------------------------------------------------------
\tikzset{
  % section band: a soft rounded region grouping cards. Draw on the
  % background layer with  fit  over the member nodes.
  tbband/.style = {
    rounded corners = 7pt, draw = tbHair, line width = 0.6pt,
    fill = tbPaper, inner sep = 6mm
  },
  tbband accent/.style = {                 % tinted band, #1 = family stem
    rounded corners = 7pt, draw = #1Line!55, line width = 0.6pt,
    fill = #1Fill!45, inner sep = 6mm
  },
  % small-caps section label, sits at a band's top-left
  tbsection/.style = {
    font = \sffamily\bfseries\scriptsize, text = tbMute, inner sep = 0pt
  },
  % figure title + subtitle (place at top, anchored north)
  tbtitle/.style = {font = \sffamily\bfseries\large, text = tbInk,
    inner sep = 0pt},
  tbsubtitle/.style = {font = \sffamily\small, text = tbMute, inner sep = 0pt},
  % center annotation / aside
  tbnote/.style = {font = \sffamily\itshape\footnotesize, text = tbMute,
    align = center, inner sep = 0pt},
  % footer principle item: icon over a short label
  tbprinciple/.style = {align = center, text = tbNeutText,
    font = \sffamily\fontsize{7}{8.5}\selectfont, inner sep = 2pt},
}

% Helper: a section label placed just above-left of a fitted band.
% \tbsectionlabel{<node name>}{<TEXT>}
\newcommand{\tbsectionlabel}[2]{%
  \node[tbsection, anchor=south west] at ([xshift=2mm,yshift=1.5mm]#1.north west)
    {\textls{\MakeUppercase{#2}}};}

% Helper: numbered step chip on a card's north-west corner.
% \tbstep{<node>}{<n>}{<family Line colour>}
\newcommand{\tbstep}[3]{\node[tbchip, fill=#3] at (#1.north west) {#2};}

% Helper: a titled card body -> icon, bold title, then small detail.
% Use inside a node:  {\tbhead{file-text}{Prospectus}\\ \tbdetail{...}}
\newcommand{\tbhead}[2]{\tbicon[4.4mm]{#1}\\[2pt]\textbf{#2}}
% Detail text: small, and with hyphenation suppressed so short labels
% wrap on word boundaries instead of breaking mid-word in narrow cards.
\newcommand{\tbdetail}[1]{{\fontsize{7}{8.6}\selectfont
  \hyphenpenalty=10000\exhyphenpenalty=10000\relax #1}}
      % path relative to the figure, or via TEXINPUTS
%   \begin{document}\begin{tikzpicture} ... \end{tikzpicture}\end{document}
%
% Design intent (read this before adding one-off styles):
%   * One palette, used SEMANTICALLY. A colour means a role, never decoration.
%   * Rounded "card" containers, thin borders, soft fills, one sans-serif.
%   * Real Lucide icons (built to PDF by icons/build_icons.py), one family.
%   * Manual layout. We choose spacing and alignment; we do not auto-route.
%
% Never hand-pick a colour in a figure. If you need a new role, add it here.
% ======================================================================

\usepackage[T1]{fontenc}
\usepackage[scaled=0.95]{helvet}      % Helvetica-like sans; swap for the
\renewcommand\familydefault{\sfdefault} % book font if the template differs.
\usepackage{graphicx}
\usepackage{amssymb}

% Letterspacing for small-caps section labels. Guarded so compilation
% still works if microtype is unavailable.
\IfFileExists{microtype.sty}{\usepackage[letterspace=140]{microtype}}{}
\providecommand{\textls}[1]{#1}   % fallback if microtype is absent

\usepackage{tikz}
\usetikzlibrary{positioning, arrows.meta, shapes.geometric, calc, fit,
  backgrounds, decorations.pathreplacing}

% ----------------------------------------------------------------------
% 1. SEMANTIC PALETTE  (keep in sync with palette.md and icons INK)
%    Each role has Fill (soft), Line (border/accent), Text (label).
%    Anchored to book-private/docs/visual-style-guide.md (the book exhibit
%    standard): ink #1F2933, muted #64748B, border #CBD5E1, light fill
%    #F8FAFC, process blue #2563EB, evidence green #059669, stop amber
%    #D97706, hard-failure red. The three CORE roles (prim / data / stop)
%    are the guide's blue / green / amber. dsgn, inf, stor are an EXTENDED
%    tier for rich conceptual maps; the guide's default is at most three
%    semantic colours per exhibit, so reach for the core three first.
% ----------------------------------------------------------------------
\definecolor{tbInk}{HTML}{1F2933}      % icons, strokes, default ink
\definecolor{tbMute}{HTML}{64748B}     % captions, section labels
\definecolor{tbLine}{HTML}{94A3B8}     % arrows (secondary stroke)
\definecolor{tbHair}{HTML}{CBD5E1}     % borders, dividers, faint panels
\definecolor{tbPaper}{HTML}{F8FAFC}    % panel / light background

% --- CORE semantic roles (the guide's three colours + neutral) ---
% neutral / slate -- the familiar, the inert, the "given"
\definecolor{tbNeutFill}{HTML}{F1F5F9}
\definecolor{tbNeutLine}{HTML}{CBD5E1}
\definecolor{tbNeutText}{HTML}{334155}
% primary / process blue -- the main path, question, project spine
\definecolor{tbPrimFill}{HTML}{EFF6FF}
\definecolor{tbPrimLine}{HTML}{2563EB}
\definecolor{tbPrimText}{HTML}{1E3A8A}
% data / evidence green -- sources, evidence, validation, checks
\definecolor{tbDataFill}{HTML}{ECFDF5}
\definecolor{tbDataLine}{HTML}{059669}
\definecolor{tbDataText}{HTML}{065F46}
% stop / amber -- stop gates, warnings, revision triggers
\definecolor{tbStopFill}{HTML}{FEF3C7}
\definecolor{tbStopLine}{HTML}{D97706}
\definecolor{tbStopText}{HTML}{78350F}

% --- EXTENDED tier (use sparingly; never break the 3-colour default
%     without a reason, e.g. a wide conceptual map) ---
% design / indigo -- design logic, comparison, identification
\definecolor{tbDsgnFill}{HTML}{EEF2FF}
\definecolor{tbDsgnLine}{HTML}{4F46E5}
\definecolor{tbDsgnText}{HTML}{3730A3}
% inference / teal -- claims, interpretation, warrants
\definecolor{tbInfFill}{HTML}{F0FDFA}
\definecolor{tbInfLine}{HTML}{0D9488}
\definecolor{tbInfText}{HTML}{115E59}
% store / violet -- the durable record, model, source of truth
\definecolor{tbStorFill}{HTML}{F5F3FF}
\definecolor{tbStorLine}{HTML}{7C3AED}
\definecolor{tbStorText}{HTML}{5B21B6}
% fail / red -- genuine hard failures only (the guide reserves red)
\definecolor{tbFailFill}{HTML}{FEF2F2}
\definecolor{tbFailLine}{HTML}{DC2626}
\definecolor{tbFailText}{HTML}{991B1B}
% backward-compatible alias: tbWarn* -> amber stop family
\colorlet{tbWarnFill}{tbStopFill}
\colorlet{tbWarnLine}{tbStopLine}
\colorlet{tbWarnText}{tbStopText}

% ----------------------------------------------------------------------
% 2. ICONS
%    Build first:  cd icons && python3 build_icons.py
%    Then: \tbicon{file-text}  or  \tbicon[5mm]{database}
%    The icon inherits the recoloured ink baked in at build time.
% ----------------------------------------------------------------------
\graphicspath{{icons/cache/}{./icons/cache/}{../icons/cache/}%
  {../../icons/cache/}{./}}
\newcommand{\tbicon}[2][4mm]{%
  \raisebox{-0.12\height}{\includegraphics[height=#1]{#2}}}

% ----------------------------------------------------------------------
% 3. NODE STYLES
% ----------------------------------------------------------------------
\tikzset{
  % spacing defaults; override per figure with node distance=...
  node distance = 9mm and 13mm,
  %
  % -- base card. Use a semantic family with  tbcard=tbData  etc.
  %    #1 family stem: Neut Prim Data Stop (core) | Dsgn Inf Stor Fail (ext).
  %    Border 0.7pt (secondary stroke, per the guide's 0.75-1.0pt band).
  tbcard/.style 2 args = {
    rounded corners = 4pt,
    draw = #1Line, line width = 0.7pt,
    fill = #1Fill, text = #1Text,
    align = #2, inner sep = 6pt,
    minimum height = 11mm, minimum width = 26mm,
    font = \sffamily\small
  },
  tbcard/.default = {tbNeut}{center},
  % convenience: one-argument families, centered text
  neut/.style = {tbcard={tbNeut}{center}},
  prim/.style = {tbcard={tbPrim}{center}},
  data/.style = {tbcard={tbData}{center}},
  stop/.style = {tbcard={tbStop}{center}},
  dsgn/.style = {tbcard={tbDsgn}{center}},
  inf/.style  = {tbcard={tbInf}{center}},
  stor/.style = {tbcard={tbStor}{center}},
  fail/.style = {tbcard={tbFail}{center}},
  warn/.style = {tbcard={tbStop}{center}},   % alias -> stop (amber)
  % left-aligned body variants (for cards with a title + detail lines)
  neutL/.style = {tbcard={tbNeut}{left}},
  primL/.style = {tbcard={tbPrim}{left}},
  dataL/.style = {tbcard={tbData}{left}},
  stopL/.style = {tbcard={tbStop}{left}},
  dsgnL/.style = {tbcard={tbDsgn}{left}},
  infL/.style  = {tbcard={tbInf}{left}},
  storL/.style = {tbcard={tbStor}{left}},
  failL/.style = {tbcard={tbFail}{left}},
  %
  % -- circular variable node for causal DAGs. Single family argument:
  %    tbvar=tbPrim (treatment), tbvar=tbInf (outcome), etc.
  tbvar/.style = {
    circle, draw = #1Line, fill = #1Fill, text = #1Text,
    line width = 0.7pt, align = center, inner sep = 1.5pt,
    minimum size = 15mm, font = \sffamily\small
  },
  tbvar/.default = tbNeut,
  %
  % -- emphasis: a heavier border for the "source of truth" / final node
  tbstrong/.style = {line width = 1.5pt},
  %
  % -- code / monospace inset (e.g. a record or schema snippet)
  tbcode/.style = {
    rounded corners = 3pt, draw = tbHair, line width = 0.5pt,
    fill = white, text = tbNeutText, align = left,
    inner sep = 5pt, font = \ttfamily\scriptsize
  },
  %
  % -- step chip: small numbered circle pinned to a card corner
  tbchip/.style = {
    circle, text = white, font = \sffamily\tiny\bfseries,
    inner sep = 0pt, minimum size = 4.2mm, fill = tbLine
  },
  %
  % -- pill tag (status label like "required" / "gated path")
  tbpill/.style = {
    rounded corners = 5pt, draw = #1Line, line width = 0.5pt,
    fill = #1Fill, text = #1Text, inner xsep = 4pt, inner ysep = 1.5pt,
    font = \sffamily\fontsize{6.5}{7.5}\selectfont
  },
  tbpill/.default = tbNeut,
}

% ----------------------------------------------------------------------
% 4. EDGES
% ----------------------------------------------------------------------
% Line-weight convention follows the guide: main path 1.25-1.5pt, secondary
% 0.75-1.0pt; solid = primary sequence / dependency, dashed = feedback /
% checking / optional. Heavier stroke marks the main path.
\tikzset{
  tbedge/.style = {-{Stealth[length=2.6mm, width=2.1mm]},
    draw = tbLine, line width = 0.9pt},               % secondary stroke
  tbedge thick/.style = {tbedge, line width = 1.4pt}, % main path
  tbdash/.style = {tbedge, dashed},                   % soft / checking link
  tbback/.style = {tbedge, dashed, draw = tbStopLine},% revision / feedback
  tbline/.style = {draw = tbLine, line width = 0.9pt},% undirected tie
}

% ----------------------------------------------------------------------
% 5. SECTIONS, TITLES, FOOTER
% ----------------------------------------------------------------------
\tikzset{
  % section band: a soft rounded region grouping cards. Draw on the
  % background layer with  fit  over the member nodes.
  tbband/.style = {
    rounded corners = 7pt, draw = tbHair, line width = 0.6pt,
    fill = tbPaper, inner sep = 6mm
  },
  tbband accent/.style = {                 % tinted band, #1 = family stem
    rounded corners = 7pt, draw = #1Line!55, line width = 0.6pt,
    fill = #1Fill!45, inner sep = 6mm
  },
  % small-caps section label, sits at a band's top-left
  tbsection/.style = {
    font = \sffamily\bfseries\scriptsize, text = tbMute, inner sep = 0pt
  },
  % figure title + subtitle (place at top, anchored north)
  tbtitle/.style = {font = \sffamily\bfseries\large, text = tbInk,
    inner sep = 0pt},
  tbsubtitle/.style = {font = \sffamily\small, text = tbMute, inner sep = 0pt},
  % center annotation / aside
  tbnote/.style = {font = \sffamily\itshape\footnotesize, text = tbMute,
    align = center, inner sep = 0pt},
  % footer principle item: icon over a short label
  tbprinciple/.style = {align = center, text = tbNeutText,
    font = \sffamily\fontsize{7}{8.5}\selectfont, inner sep = 2pt},
}

% Helper: a section label placed just above-left of a fitted band.
% \tbsectionlabel{<node name>}{<TEXT>}
\newcommand{\tbsectionlabel}[2]{%
  \node[tbsection, anchor=south west] at ([xshift=2mm,yshift=1.5mm]#1.north west)
    {\textls{\MakeUppercase{#2}}};}

% Helper: numbered step chip on a card's north-west corner.
% \tbstep{<node>}{<n>}{<family Line colour>}
\newcommand{\tbstep}[3]{\node[tbchip, fill=#3] at (#1.north west) {#2};}

% Helper: a titled card body -> icon, bold title, then small detail.
% Use inside a node:  {\tbhead{file-text}{Prospectus}\\ \tbdetail{...}}
\newcommand{\tbhead}[2]{\tbicon[4.4mm]{#1}\\[2pt]\textbf{#2}}
% Detail text: small, and with hyphenation suppressed so short labels
% wrap on word boundaries instead of breaking mid-word in narrow cards.
\newcommand{\tbdetail}[1]{{\fontsize{7}{8.6}\selectfont
  \hyphenpenalty=10000\exhyphenpenalty=10000\relax #1}}
      % path relative to the figure, or via TEXINPUTS
%   \begin{document}\begin{tikzpicture} ... \end{tikzpicture}\end{document}
%
% Design intent (read this before adding one-off styles):
%   * One palette, used SEMANTICALLY. A colour means a role, never decoration.
%   * Rounded "card" containers, thin borders, soft fills, one sans-serif.
%   * Real Lucide icons (built to PDF by icons/build_icons.py), one family.
%   * Manual layout. We choose spacing and alignment; we do not auto-route.
%
% Never hand-pick a colour in a figure. If you need a new role, add it here.
% ======================================================================

\usepackage[T1]{fontenc}
\usepackage[scaled=0.95]{helvet}      % Helvetica-like sans; swap for the
\renewcommand\familydefault{\sfdefault} % book font if the template differs.
\usepackage{graphicx}
\usepackage{amssymb}

% Letterspacing for small-caps section labels. Guarded so compilation
% still works if microtype is unavailable.
\IfFileExists{microtype.sty}{\usepackage[letterspace=140]{microtype}}{}
\providecommand{\textls}[1]{#1}   % fallback if microtype is absent

\usepackage{tikz}
\usetikzlibrary{positioning, arrows.meta, shapes.geometric, calc, fit,
  backgrounds, decorations.pathreplacing}

% ----------------------------------------------------------------------
% 1. SEMANTIC PALETTE  (keep in sync with palette.md and icons INK)
%    Each role has Fill (soft), Line (border/accent), Text (label).
%    Anchored to book-private/docs/visual-style-guide.md (the book exhibit
%    standard): ink #1F2933, muted #64748B, border #CBD5E1, light fill
%    #F8FAFC, process blue #2563EB, evidence green #059669, stop amber
%    #D97706, hard-failure red. The three CORE roles (prim / data / stop)
%    are the guide's blue / green / amber. dsgn, inf, stor are an EXTENDED
%    tier for rich conceptual maps; the guide's default is at most three
%    semantic colours per exhibit, so reach for the core three first.
% ----------------------------------------------------------------------
\definecolor{tbInk}{HTML}{1F2933}      % icons, strokes, default ink
\definecolor{tbMute}{HTML}{64748B}     % captions, section labels
\definecolor{tbLine}{HTML}{94A3B8}     % arrows (secondary stroke)
\definecolor{tbHair}{HTML}{CBD5E1}     % borders, dividers, faint panels
\definecolor{tbPaper}{HTML}{F8FAFC}    % panel / light background

% --- CORE semantic roles (the guide's three colours + neutral) ---
% neutral / slate -- the familiar, the inert, the "given"
\definecolor{tbNeutFill}{HTML}{F1F5F9}
\definecolor{tbNeutLine}{HTML}{CBD5E1}
\definecolor{tbNeutText}{HTML}{334155}
% primary / process blue -- the main path, question, project spine
\definecolor{tbPrimFill}{HTML}{EFF6FF}
\definecolor{tbPrimLine}{HTML}{2563EB}
\definecolor{tbPrimText}{HTML}{1E3A8A}
% data / evidence green -- sources, evidence, validation, checks
\definecolor{tbDataFill}{HTML}{ECFDF5}
\definecolor{tbDataLine}{HTML}{059669}
\definecolor{tbDataText}{HTML}{065F46}
% stop / amber -- stop gates, warnings, revision triggers
\definecolor{tbStopFill}{HTML}{FEF3C7}
\definecolor{tbStopLine}{HTML}{D97706}
\definecolor{tbStopText}{HTML}{78350F}

% --- EXTENDED tier (use sparingly; never break the 3-colour default
%     without a reason, e.g. a wide conceptual map) ---
% design / indigo -- design logic, comparison, identification
\definecolor{tbDsgnFill}{HTML}{EEF2FF}
\definecolor{tbDsgnLine}{HTML}{4F46E5}
\definecolor{tbDsgnText}{HTML}{3730A3}
% inference / teal -- claims, interpretation, warrants
\definecolor{tbInfFill}{HTML}{F0FDFA}
\definecolor{tbInfLine}{HTML}{0D9488}
\definecolor{tbInfText}{HTML}{115E59}
% store / violet -- the durable record, model, source of truth
\definecolor{tbStorFill}{HTML}{F5F3FF}
\definecolor{tbStorLine}{HTML}{7C3AED}
\definecolor{tbStorText}{HTML}{5B21B6}
% fail / red -- genuine hard failures only (the guide reserves red)
\definecolor{tbFailFill}{HTML}{FEF2F2}
\definecolor{tbFailLine}{HTML}{DC2626}
\definecolor{tbFailText}{HTML}{991B1B}
% backward-compatible alias: tbWarn* -> amber stop family
\colorlet{tbWarnFill}{tbStopFill}
\colorlet{tbWarnLine}{tbStopLine}
\colorlet{tbWarnText}{tbStopText}

% ----------------------------------------------------------------------
% 2. ICONS
%    Build first:  cd icons && python3 build_icons.py
%    Then: \tbicon{file-text}  or  \tbicon[5mm]{database}
%    The icon inherits the recoloured ink baked in at build time.
% ----------------------------------------------------------------------
\graphicspath{{icons/cache/}{./icons/cache/}{../icons/cache/}%
  {../../icons/cache/}{./}}
\newcommand{\tbicon}[2][4mm]{%
  \raisebox{-0.12\height}{\includegraphics[height=#1]{#2}}}

% ----------------------------------------------------------------------
% 3. NODE STYLES
% ----------------------------------------------------------------------
\tikzset{
  % spacing defaults; override per figure with node distance=...
  node distance = 9mm and 13mm,
  %
  % -- base card. Use a semantic family with  tbcard=tbData  etc.
  %    #1 family stem: Neut Prim Data Stop (core) | Dsgn Inf Stor Fail (ext).
  %    Border 0.7pt (secondary stroke, per the guide's 0.75-1.0pt band).
  tbcard/.style 2 args = {
    rounded corners = 4pt,
    draw = #1Line, line width = 0.7pt,
    fill = #1Fill, text = #1Text,
    align = #2, inner sep = 6pt,
    minimum height = 11mm, minimum width = 26mm,
    font = \sffamily\small
  },
  tbcard/.default = {tbNeut}{center},
  % convenience: one-argument families, centered text
  neut/.style = {tbcard={tbNeut}{center}},
  prim/.style = {tbcard={tbPrim}{center}},
  data/.style = {tbcard={tbData}{center}},
  stop/.style = {tbcard={tbStop}{center}},
  dsgn/.style = {tbcard={tbDsgn}{center}},
  inf/.style  = {tbcard={tbInf}{center}},
  stor/.style = {tbcard={tbStor}{center}},
  fail/.style = {tbcard={tbFail}{center}},
  warn/.style = {tbcard={tbStop}{center}},   % alias -> stop (amber)
  % left-aligned body variants (for cards with a title + detail lines)
  neutL/.style = {tbcard={tbNeut}{left}},
  primL/.style = {tbcard={tbPrim}{left}},
  dataL/.style = {tbcard={tbData}{left}},
  stopL/.style = {tbcard={tbStop}{left}},
  dsgnL/.style = {tbcard={tbDsgn}{left}},
  infL/.style  = {tbcard={tbInf}{left}},
  storL/.style = {tbcard={tbStor}{left}},
  failL/.style = {tbcard={tbFail}{left}},
  %
  % -- circular variable node for causal DAGs. Single family argument:
  %    tbvar=tbPrim (treatment), tbvar=tbInf (outcome), etc.
  tbvar/.style = {
    circle, draw = #1Line, fill = #1Fill, text = #1Text,
    line width = 0.7pt, align = center, inner sep = 1.5pt,
    minimum size = 15mm, font = \sffamily\small
  },
  tbvar/.default = tbNeut,
  %
  % -- emphasis: a heavier border for the "source of truth" / final node
  tbstrong/.style = {line width = 1.5pt},
  %
  % -- code / monospace inset (e.g. a record or schema snippet)
  tbcode/.style = {
    rounded corners = 3pt, draw = tbHair, line width = 0.5pt,
    fill = white, text = tbNeutText, align = left,
    inner sep = 5pt, font = \ttfamily\scriptsize
  },
  %
  % -- step chip: small numbered circle pinned to a card corner
  tbchip/.style = {
    circle, text = white, font = \sffamily\tiny\bfseries,
    inner sep = 0pt, minimum size = 4.2mm, fill = tbLine
  },
  %
  % -- pill tag (status label like "required" / "gated path")
  tbpill/.style = {
    rounded corners = 5pt, draw = #1Line, line width = 0.5pt,
    fill = #1Fill, text = #1Text, inner xsep = 4pt, inner ysep = 1.5pt,
    font = \sffamily\fontsize{6.5}{7.5}\selectfont
  },
  tbpill/.default = tbNeut,
}

% ----------------------------------------------------------------------
% 4. EDGES
% ----------------------------------------------------------------------
% Line-weight convention follows the guide: main path 1.25-1.5pt, secondary
% 0.75-1.0pt; solid = primary sequence / dependency, dashed = feedback /
% checking / optional. Heavier stroke marks the main path.
\tikzset{
  tbedge/.style = {-{Stealth[length=2.6mm, width=2.1mm]},
    draw = tbLine, line width = 0.9pt},               % secondary stroke
  tbedge thick/.style = {tbedge, line width = 1.4pt}, % main path
  tbdash/.style = {tbedge, dashed},                   % soft / checking link
  tbback/.style = {tbedge, dashed, draw = tbStopLine},% revision / feedback
  tbline/.style = {draw = tbLine, line width = 0.9pt},% undirected tie
}

% ----------------------------------------------------------------------
% 5. SECTIONS, TITLES, FOOTER
% ----------------------------------------------------------------------
\tikzset{
  % section band: a soft rounded region grouping cards. Draw on the
  % background layer with  fit  over the member nodes.
  tbband/.style = {
    rounded corners = 7pt, draw = tbHair, line width = 0.6pt,
    fill = tbPaper, inner sep = 6mm
  },
  tbband accent/.style = {                 % tinted band, #1 = family stem
    rounded corners = 7pt, draw = #1Line!55, line width = 0.6pt,
    fill = #1Fill!45, inner sep = 6mm
  },
  % small-caps section label, sits at a band's top-left
  tbsection/.style = {
    font = \sffamily\bfseries\scriptsize, text = tbMute, inner sep = 0pt
  },
  % figure title + subtitle (place at top, anchored north)
  tbtitle/.style = {font = \sffamily\bfseries\large, text = tbInk,
    inner sep = 0pt},
  tbsubtitle/.style = {font = \sffamily\small, text = tbMute, inner sep = 0pt},
  % center annotation / aside
  tbnote/.style = {font = \sffamily\itshape\footnotesize, text = tbMute,
    align = center, inner sep = 0pt},
  % footer principle item: icon over a short label
  tbprinciple/.style = {align = center, text = tbNeutText,
    font = \sffamily\fontsize{7}{8.5}\selectfont, inner sep = 2pt},
}

% Helper: a section label placed just above-left of a fitted band.
% \tbsectionlabel{<node name>}{<TEXT>}
\newcommand{\tbsectionlabel}[2]{%
  \node[tbsection, anchor=south west] at ([xshift=2mm,yshift=1.5mm]#1.north west)
    {\textls{\MakeUppercase{#2}}};}

% Helper: numbered step chip on a card's north-west corner.
% \tbstep{<node>}{<n>}{<family Line colour>}
\newcommand{\tbstep}[3]{\node[tbchip, fill=#3] at (#1.north west) {#2};}

% Helper: a titled card body -> icon, bold title, then small detail.
% Use inside a node:  {\tbhead{file-text}{Prospectus}\\ \tbdetail{...}}
\newcommand{\tbhead}[2]{\tbicon[4.4mm]{#1}\\[2pt]\textbf{#2}}
% Detail text: small, and with hyphenation suppressed so short labels
% wrap on word boundaries instead of breaking mid-word in narrow cards.
\newcommand{\tbdetail}[1]{{\fontsize{7}{8.6}\selectfont
  \hyphenpenalty=10000\exhyphenpenalty=10000\relax #1}}


\begin{document}
\begin{tikzpicture}[
  layer/.style = {tbcard={#1}{center}, text width=52mm,
    minimum width=60mm, minimum height=15mm, inner sep=5pt,
    font=\sffamily\fontsize{8.5}{10.5}\selectfont,
    execute at begin node={\hyphenpenalty=10000\exhyphenpenalty=10000}},
  elabel/.style = {font=\sffamily\fontsize{7.5}{9}\selectfont,
    text=tbMute, fill=white, inner sep=1.5pt, align=center},
]

% ------------------------------------------------ the downward spine
\node[layer={tbNeut}] (you) at (0, 0)
  {\tbhead{keyboard}{You}\\[2pt]
   \tbdetail{type one line and press return}};

\node[layer={tbNeut}] (term) at (0, -30mm)
  {\tbhead{monitor}{Terminal}\\[2pt]
   \tbdetail{a window that draws text and passes
   your keystrokes on; it runs nothing itself}};

\node[layer={tbPrim}, tbstrong] (shell) at (0, -62mm)
  {\tbhead{square-terminal}{Shell (bash, zsh)}\\[2pt]
   \tbdetail{reads the line, expands variables and globs,
   finds the program on PATH, then asks the kernel
   to run it}};

\node[layer={tbNeut}] (kern) at (0, -94mm)
  {\tbhead{cog}{Operating system kernel}\\[2pt]
   \tbdetail{creates and supervises processes; owns the
   filesystem, memory, and devices}};

\node[layer={tbData}, tbstrong] (proc) at (0, -128mm)
  {\tbhead{cpu}{Your command, running: a process}\\[2pt]
   \tbdetail{PID $\cdot$ working directory $\cdot$ a copy of the
   environment $\cdot$ three streams: stdin, stdout, stderr}};

% ------------------------------------------------ downward edges
\draw[tbedge thick] (you) -- (term)
  node[elabel, midway] {keystrokes};
\draw[tbedge thick] (term) -- (shell)
  node[elabel, midway] {the line you typed};
\draw[tbedge thick] (shell) -- (kern)
  node[elabel, midway] {run this program with these arguments};
\draw[tbedge thick] (kern) -- (proc)
  node[elabel, midway] {starts};

% ------------------------------------------------ return edges
% stdout/stderr flow back up to the terminal (right side).
\draw[tbedge, draw=tbDataLine]
  (proc.east) .. controls +(22mm,0) and +(22mm,0) .. (term.east)
  node[elabel, text=tbDataText, pos=0.5, xshift=1.5mm, anchor=west]
  {stdout and stderr\\ (what you see\\ printed)};

% exit code reports back to the shell (left side, feedback dash).
\draw[tbback]
  (proc.west) .. controls +(-20mm,0) and +(-20mm,0) .. (shell.west)
  node[elabel, text=tbStopText, pos=0.5, xshift=-1.5mm, anchor=east]
  {exit code\\ (0 = success)};

\end{tikzpicture}
\end{document}
